\documentclass[12pt]{article}

\usepackage[margin=1.12in]{geometry}
\usepackage{amsmath,amssymb,amsthm,mathtools}
\usepackage{booktabs}
\usepackage{hyperref}

\newtheorem{theorem}{Theorem}
\newtheorem{proposition}{Proposition}
\theoremstyle{remark}
\newtheorem{remark}{Remark}

\newcommand{\B}{\mathcal B}
\newcommand{\K}{\mathcal K}
\newcommand{\M}{\mathcal M}
\newcommand{\R}{\mathbb R}
\newcommand{\id}{\operatorname{id}}

\setlength{\emergencystretch}{3em}
\hypersetup{
  pdftitle={Crossed-Product Entropy Does Not Supply an Omitted Finite-Center Calibration},
  pdfsubject={Type II entropy, modular crossed products, finite centers, black-hole entropy},
  pdfkeywords={crossed product, Type II algebra, entropy calibration, finite centers, black holes}
}

\title{Crossed-Product Entropy Does Not Supply\\
an Omitted Finite-Center Calibration}
\author{}
\date{Live review note of June 20, 2026}

\begin{document}
\maketitle

\begin{abstract}
Gravitational crossed-product constructions define entropy for a specified
Type II or semifinite von Neumann algebra together with its trace or
normalization convention.  This note isolates a finite-central boundary.  If
\(\M\) is replaced by
\[
  \B_n=\M\,\bar\otimes\,\ell^\infty(\{1,\ldots,n\}),
\]
then the modular crossed product carries the finite center to the Type II core:
\[
  \B_n\rtimes_{\sigma^{\Omega_r}}\R
  \cong
  (\M\rtimes_{\sigma^\omega}\R)\,\bar\otimes\,
  \ell^\infty(\{1,\ldots,n\}).
\]
Trace-normalized Type II entropy fixes the finite-center entropy contribution
only when the supplied trace or absolute calibration separates the central
weights relevant to the claimed state family.  Therefore Witten-style
crossed-product entropy is a positive entropy construction for the supplied
algebra and trace; it is not, by itself, a principle that reconstructs an
omitted finite central summand, finite central trace weights, or an absolute
entropy-density coefficient.  The positive exit is equally concrete: atomic
trace calibration on the retained finite central summands recovers the
finite-center entropy term up to the intended common trace convention.
\end{abstract}

\paragraph{Status.}
This is non-frozen external-review metadata for Team 1.  It is intended as a
small claim that a crossed-product, Type II, or black-hole entropy reviewer can
classify as known, false, missing-hypothesis, or useful.  It does not modify
any frozen Team 1 manuscript, Lean archive, public mirror manifest, or SHA-256
package.

\section{Primary sources}

The boundary below was checked against Witten's crossed-product black-hole
entropy construction \cite{WittenCrossed}, Longo--Witten continuous entropy
\cite{LongoWittenEntropy}, the large-\(N\) generalized-entropy algebra of
Chandrasekaran--Penington--Witten \cite{CPWLargeN}, and the JT-gravity Type II
algebra analysis of Penington--Witten \cite{PeningtonWittenJT}.  These works
are used here as positive constructions for specified algebras and traces, not
as claims that finite-center calibration is absent from physics.

\section{Finite centers in the crossed product}

Let \(\M\) be a von Neumann algebra with faithful normal state \(\omega\), and
let
\[
  \K=\M\rtimes_{\sigma^\omega}\R
\]
be its modular crossed product.  For \(n\ge2\), let
\[
  \B_n=\M\,\bar\otimes\,\ell^\infty(\{1,\ldots,n\}),
  \qquad
  \Omega_r=\omega\otimes r,
\]
where \(r=(r_1,\ldots,r_n)\) is a faithful probability vector.

\begin{theorem}[finite central cores]
\label{thm:finite-central-cores}
There is a canonical crossed-product identification
\[
  \B_n\rtimes_{\sigma^{\Omega_r}}\R
  \cong
  \K\,\bar\otimes\,\ell^\infty(\{1,\ldots,n\}).
\]
Consequently the modular crossed product carries the finite center into the
Type II core.  It does not by itself select the finite central cardinality,
the predicate \(n=2\), or the finite-center trace normalization.
\end{theorem}

\begin{proof}
The modular automorphism group of \(\Omega_r=\omega\otimes r\) is
\[
  \sigma_t^{\Omega_r}=\sigma_t^\omega\otimes\id
\]
because the finite abelian factor is central and \(r\) is diagonal on it.
Crossed products commute with tensoring by a finite algebra on which the
action is trivial, giving the displayed isomorphism.  The right-hand side
still has central projections \(1\otimes e_i\), so the finite center has not
been selected or removed by the passage to the core.
\end{proof}

\section{Trace calibration fiber}

Assume \(\K\) is equipped with a faithful normal semifinite trace \(\tau\).
On the full finite central core, consider traces
\[
  \tau_a(x_1,\ldots,x_n)=\sum_{i=1}^n a_i\tau(x_i),
  \qquad a_i>0.
\]
For a central probability vector \(p=(p_i)\), the part of the finite-center
entropy that changes with the central trace weights is
\[
  L_a(p)=\sum_i p_i\log a_i.
\]
If the supplied trace calibration consists of central test values
\[
  T_a(c)=\sum_i a_i c_i,\qquad c\in C\subset\R^n,
\]
then \(L_a(p)\) is determined for all allowed \(p\in P\) if and only if
\[
  T_a|_C=T_b|_C
  \quad\Longrightarrow\quad
  L_a(p)=L_b(p)\quad\hbox{for every }p\in P .
  \tag{1}
\]

\begin{proposition}[diagonal trace tests are insufficient]
\label{prop:diagonal-insufficient}
Scalar diagonal trace tests do not determine the finite-center entropy
normalization.  For \(n=2\), the trace weights \(a=(1,3)\) and \(b=(2,2)\)
agree on every scalar diagonal test, but for \(p=(1/2,1/2)\),
\[
  L_a(p)-L_b(p)=\frac12\log\frac34\ne0 .
\]
\end{proposition}

\begin{proof}
A scalar diagonal test has the form \(c=(\lambda,\lambda)\), so
\[
  T_a(c)=\lambda(a_1+a_2)=4\lambda=T_b(c).
\]
However
\[
  L_a(1/2,1/2)-L_b(1/2,1/2)
  =
  \frac12(\log1+\log3-\log2-\log2)
  =
  \frac12\log\frac34 .
\]
Thus the same diagonal Type II trace data can leave different absolute
finite-center entropy terms.
\end{proof}

\begin{proposition}[atomic calibration recovers the finite-center entropy term]
\label{prop:atomic-calibration}
Assume the full finite central core
\(\K\,\bar\otimes\,\ell^\infty(\{1,\ldots,n\})\) is retained and the supplied
calibration contains, for some positive \(h\in\K\) with \(0<\tau(h)<\infty\),
the atomic trace values
\[
  \tau_a(h\otimes e_i)=a_i\tau(h),\qquad i=1,\ldots,n .
\]
Then the central trace weights \(a_i\) are determined up to the same convention
used to normalize \(\tau(h)\).  In particular, for every retained central
probability vector \(p\), the finite-center entropy contribution
\[
  L_a(p)=\sum_i p_i\log a_i
\]
is determined.  If only the projective class \([a]\) is fixed, then \(L_a(p)\)
is determined up to a common additive constant independent of \(p\), and all
differences \(L_a(p)-L_a(p')\) are determined exactly.
\end{proposition}

\begin{proof}
Dividing the supplied value \(\tau_a(h\otimes e_i)\) by the fixed normalization
\(\tau(h)\) gives \(a_i\).  Substitution gives \(L_a(p)\).  If the trace is
specified only up to common rescaling \(a\mapsto\lambda a\), then
\[
  L_{\lambda a}(p)=\sum_i p_i\log(\lambda a_i)
  =L_a(p)+\log\lambda,
\]
because \(\sum_i p_i=1\).  The shift is independent of \(p\), so it cancels in
differences between retained central states.
\end{proof}

\section{Entropy-density boundary}

Let \(D\) be the crossed-product, relative-entropy, or operational data
supplied in a comparison class, and let
\[
  c(X)=\lim_\Sigma\frac{S_X(\Sigma)}{\operatorname{Area}(\Sigma)}
\]
be an asserted absolute entropy-density coefficient.  Then \(c\) is determined
by \(D\) exactly when it is constant on \(D\)-fibers:
\[
  D(X)=D(Y)
  \quad\Longrightarrow\quad
  \lim_\Sigma
  \frac{S_X(\Sigma)-S_Y(\Sigma)}
       {\operatorname{Area}(\Sigma)}
  =0 .
  \tag{2}
\]
Finite-center cardinality terms, trace-weight terms such as \(L_a-L_b\), and
affine entropy or area counterterms are obstructions precisely when their
normalized limits in (2) are nonzero.

\section{Positive exits and referee question}

\begin{center}
\begin{tabular}{p{0.35\linewidth}p{0.50\linewidth}}
\toprule
Exit premise & Effect \\
\midrule
Full Type II algebra retained & The finite central summands are part of the
specified entropy algebra, rather than omitted before the entropy calculation.
\\
\addlinespace
Atomic trace calibration & Tests \(h\otimes e_i\) recover the central trace
weights \(a_i\), hence the finite-center entropy contribution, by
Proposition~\ref{prop:atomic-calibration}. \\
\addlinespace
Absolute area/entropy calibration & An affinely spanning calibrated state
family fixes the additive or affine entropy residual. \\
\addlinespace
Quotient, factor, or inadmissibility axiom & The finite-center comparison
class is explicitly removed from the target theory. \\
\bottomrule
\end{tabular}
\end{center}

The precise external question is:

\begin{quote}
Does the crossed-product or gravitational Type II construction being invoked
in the target setting supply the finite central trace/area calibration above,
or does it define entropy only for the algebra and trace normalization already
specified?
\end{quote}

If the answer is the second, Team 1 has no new crossed-product theorem.  Its
surviving claim is only the claim-boundary warning: crossed-product entropy
should not be overread as finite-center maximality or as an absolute
finite-center calibration principle.

\begin{thebibliography}{9}

\bibitem{WittenCrossed}
E. Witten,
\emph{Gravity and the Crossed Product},
JHEP 10 (2022), 008,
\href{https://arxiv.org/abs/2112.12828}{arXiv:2112.12828}.

\bibitem{LongoWittenEntropy}
R. Longo and E. Witten,
\emph{A note on continuous entropy},
\href{https://arxiv.org/abs/2202.03357}{arXiv:2202.03357}.

\bibitem{CPWLargeN}
V. Chandrasekaran, G. Penington, and E. Witten,
\emph{Large N algebras and generalized entropy},
JHEP 04 (2023), 009,
\href{https://arxiv.org/abs/2209.10454}{arXiv:2209.10454}.

\bibitem{PeningtonWittenJT}
G. Penington and E. Witten,
\emph{Algebras and States in JT Gravity},
\href{https://arxiv.org/abs/2301.07257}{arXiv:2301.07257}.

\end{thebibliography}

\end{document}
